\chapter{Il ricorso}
Ed eccoci a casa di Laura. \'E arrivata in tempo a predere la sorella e ora Valentina sta finalmente sul tappeto  con le sue matite e i suoi pennarelli. Laura invece  è al lavello che prepara la cena con  un orecchio  teso cerca di decifrare i  borbottii e i sussurri che riesce a captare. Mezze frasi, accenni  a mondi fantastici. Continenti a forma di frutta, paesi costellati da piccoli laghetti, villaggi fatti solo da alberi,  terre aride e disastrate dove vivono popolazioni tornate alla barbarie. Tutto in un foglio di una cinquantina di centimetri per quaranta.

Poi, quando l'odore delle zucchine  raggiunge il naso della piccola,  un pensiero più concreto sospende tempraneamente la sua attività, lasciando la priorità alla contingenza:
«Cosa prepari per cena?» domanda.\\
«Tu cosa avresti voglia di mangiare?»

\par\medskip
\begin{center}
  \includegraphics[width=.7\linewidth]{soggiorno_tiny.png}%
\end{center}

«I bastoncini fritti! Prepari i bastoncini fritti?»\\
«Infatti.»\\
«Davvero ci sono i bastoncini?»\\
«No, c’è il riso integrale con le verdure.»\\
«Perché sempre le verdure?»\\
«Perché ce ne abbiamo nell’orto, sono buone, fanno bene e non escono dal tubo.»\\
«Blaa...il pastone...» Valentina fa una smorfia. «Ma io ho voglia di bastoncini. Quando li mangiamo?»\\

«Venerdì. A fine settimana acquisterò del cibo. Dobbiamo mangiare tutti i principi nutritivi.»

Da sotto il tavolo Rocky si alza e si avvicina alla porta. Laura lo segue con lo sguardo «Cosa hai visto?»\\
Rocky abbaia. Laura si ferma di tagliare: «Guarda, arriva Caterina. Non farle vedere che fai i capricci per la cena.»\\
«Non è vero che faccio i capricci.» Intanto l'amica è arrivata alla porta, bussa e fa scorrere l'anta di vetro, ma si arresta sull'uscio.\\
«Ciao Caterina, vieni dentro» la invita Valentina.\\

«Ma guarda che confusione!» la voce di Caterina ha qualche accento di rimprovero. «Penne, pennarelli» continua,  «lapis, fogli da disegno mi sbarrano la strada e un Cerbero feroce non mi lascia passare! Come posso entrare, signorina?»\\
«Salta!» le risponde, «E tu seduto, Rocky!».



\par\medskip
\begin{center}
  \includegraphics[width=1\linewidth]{soggiorno_prospettiva.png}%
\end{center}
Si alza dal tappeto, prende Caterina per mano, la trascina dentro con un piccolo salto sopra la distesa mentre Rocky torna a sdraiarsi sotto il tavolo.
Valentina lo guarda: «Non ha tre teste...» protesta.\\
Il grande foglio che hanno appena saltato attira l'attenzione dell'amica che vi si inginocchia vicino.

«Cosa disegni? \'E un compito?»\\
Valentina scuote veloce la testa da una parte all'altra come una campana: «Anche, ma non solo. Stavo disegnando una cosa. Vedi questo mare?»\\
Caterina si avvicina: «Si, le vedo» conferma.\\
«Divide due regni. Da una parte ci stanno i buoni»  ed indica il disegno di due uomini visibilmente forti, «dall'altra loro, che sono i cattivi.» Questa volta sposta il dito su tre figure maschili con tratti feroci.


«Fanno proprio paura» conferma Caterina cercando Laura con lo sguardo. Lei però è ancora intenta alla preparazione della cena, ma colta una vena di preoccupazione nella voce dell'amica la tranquillizza: «Fanno paura anche a me,» le conferma, «ma ha trovato la libreria manga dei miei e adesso non la ferma più nessuno.» \\
Caterina non risponde, Valentina l'ha preso un braccio tirandosela vicina. «Ti fanno paura? Davvero?»\\
«Il tuo disegno è bello, ma loro sono davvero brutti. Quindi?» Le chiede, «In che regno sarebbero questi brutti mostri?»\\
Valentina appoggia il foglio, «Loro abitano nella terra degli Shura.»\\
«E dove si trova?» le chiede.\\
Lo riprende in mano  e indica la regione dove abitano gli eroi buoni: «Se questo è il giappone, questa cosa sarà mai?»\\
Le orecchie di Laura si tendono verso la sorella «Non fare la maestrina con Caterina.» Valentina la ignora e continua: «Dai, che paese si trova a sinistra del Giappone?»\\
«Allora... abbiamo detto che questo è il Giappone? Direi che questa dovrebbe essere» attende un attimio, «dovrebbe essere la Corea» esclama, sicura della correttezza della risposta, ma la delusione non tarda ad arrivare.\\
«Hai studiato geografia da piccola? Lì c'è la Cina» e ciò detto indica con il dito la vicinanza tra le due terre.
«Hai ragione Vale, la Corea è un po' più in là.»
Caterina si inginocchia, pronta ad ascoltare la storia di quei luoghi che sicuramente non è ancora conclusa. Valentina infatti non sembra proprio aver esaurito gli argomenti:

«Ecco, vedi, qui in Giappone invece sono buoni. Loro studiano sempre le arti marziali, in questa scuola» e indica una casetta.
Caterina  avvicina gli occhi al disegno.

«Sì? E in questa scuola c’è anche un maestro?»\\
«Certo, è il padre dei tre fratelli. Adesso ti spiego bene» e parlando si posiziona veloce un cuscino sotto le ginocchia, ma viene interrotta dalla sorella:
«Vale, forse Caterina doveva dirmi qualcosa. Sentiamo?»\\
Caterina sorride,
«Si, ma lascia prima che mi racconti. È passato tanto tempo da quando Alice  mi mostrava i suoi disegni. Non pensavo mi mancasse così tanto.»
\\
«Anche Alice fa i disegni? La prossima volta voglio vederli.»\\

«Li faceva da piccola tesoro, adesso mi scrive le canzoni.»
\\
«Che bello, allora le voglio ascoltare!»\\

Caterina le sorride e appoggia il disegno sul tavolo seguita dallo sguardo corruciato di Valentina che prontamente lo va a riprendere.\\

«\'E andata male» il tono è quello riservato a Laura, non parla più con Valentina. «Mi hanno negato il visto. Ho tempo fino alle ventiquattro per fare ricorso.»\\

Laura appoggia il coltello sul tagliere e guarda verso l’amica:

«Quando l’hai saputo?»\\
«Due ore fa eri a casa mia e non lo sapevo ancora... quindi?»\\
«Hai ragione, Cate, scusami.»\\
«No, scusami tu. Sono stata acida. È che sono disperata.»\\
«Ma hai già sentito Mark?»\\
«Sì, è stato lui a dirmi della possibilità di fare ricorso.»\\
«Oh, Laura, come facciamo?»

Laura si avvicina all'amica e la prende delicatamente per le braccia:

«Cate, ne abbiamo passate di peggio. Se il tempo stringe, vediamo prima il ricorso, poi mangiamo.»\\

Valentina le ascolta distrattamente:
«Non ceniamo?» chiede, ma non ricevendo risposta dimentica la fame come solo i bambini sanno fare e ritorna a colorare.\\

In camera da letto ha un tavolino basso, lo prende insieme al suo computer  portatile e lo appoggia tra il tappeto e il frigo: «Mi fai un po' di posto Vale?»\\
«Cosa fate con il computer? Posso vedere?»\\
«Cose da grandi» le risponde Laura. Valentina non replica, si siede a fianco di Caterina. Laura collega il dispositivo ad una rete dati separata da quella domestica:

«Vediamo, quindi Mark lavora per la SAGA-M, giusto?» La connessione è stabilita e Laura accede a un servizio di informazioni basato su un protocollo dati testuale, un po' datato ma ancora in uso. È alla ricerca della SAGA-M. Caterina allunga il collo verso lo schermo:
«Sì, perché? Cosa stai cercando?»\\

«Voglio capire cosa le interessa e cosa invece la preoccupa...»\\
«Laura ti prego sii più chiara... mi stai facendo salire l'ansia, non guardiamo la procedura per il ricorso?»\\
Laura chiude alcune schede poi avvia un monitor a caratteri, uno sguardo rapido all'amica: «Scusami... hai ragione. Adesso... vedi, tu non sei mai stata negli Stati Uniti, giusto?»\\
«No, è...» si ferma, «sarebbe la prima volta.» Risponde Caterina. Valentina si sdraia sul tappeto, prende un pennarello e si disegna un cuore sulla mano.\\
«Quindi non hai commesso nulla là che possa precludere il tuo visto... Però è possibile che un'informazione su di te attivi un qualche flag di allarme...»\\
«Ma la privacy?»\\
«...La privacy riguarda la tua sfera privata, non quello che fai in pubblico...»\\
«Non capisco bene il ragionamento Laura...» L'attenzione per un attimo è attratta da Valentina, si è tolta le calze e si sta colorando anche i piedi, Caterina non vorrebbe, ma le scappa da ridere. Si riprende, è preoccupata per il suo viaggio: «Ti servono informazioni per fare ricorso?»\\
«La possibilità di fare ricorso è più apparenza che sostanza, nella pratica è un algoritmo che decide e se non capiamo davvero cosa abbia scatenato il rifiuto...»
Laura marca alcune schede dati per l'elaborazione, «ecco la scheda della compagnia...»\\
\begin{terminalbox}
EU Transparency Register
Entry extract / Interests represented
Reference: EUTR/SAGA-M/ENER/2041-06
Registrant: SAGA-M European Affairs Office

Settore: energia, reti intelligenti,
resilienza urbana.
Policy UE: Green Deal, AI Act,
infrastrutture energetiche.
Comunicazione: autonomia abitativa, mobilita' individuale,
riduzione della dipendenza da nuclei
familiari stabili.

Campagne:
- Single by design
- Non devi nulla a nessuno
- La cura e' un servizio, non un legame
\end{terminalbox}

«Quindi?»\\
Laura inspira profondamente, guarda l'amica:
«Cosa c'entra il "Single by Design" con il green deal?»\\
«Laura mi stai confondendo...»\\

La fissa per un istante, poi abbassa lo sguardo sullo schermo.\\

«Anche io sono confusa, ma se non hai commesso reati e non sei mai stata là, il tuo visto è compromesso dalla tua relazione con Mark e dal suo lavoro... Almeno... Comunque adesso vediamo la procedura per il ricorso, ma poi dobbiamo individuare il nodo che ti ha bloccata, altrimenti il ricorso sarà inutile.»

